% 04-method-2026 — what changed computationally since 1990
\section{The numerical method}
\label{sec:method}

The 1990 program generated all colour-singlet partitions of $K$, removed the
overcomplete directions by diagonalizing the inner-product matrix, and
diagonalized the Hamiltonian in the surviving basis --- storing the free and
interacting parts separately so that a change of coupling is a rescaling
rather than a rebuild~\cite{Hornbostel:1988fb}.  Nothing in that architecture
needed changing; what changed is what it is built on.  We run two
implementations against each other: the original Fortran, kept deliberately
unmodified as a historical reference, and an open-source Python port whose
matrix builds are verified \emph{bit-identical} to their own independent
reference path, with physics agreement between the two solvers at the level
of the arithmetic rather than the plot (Appendix~\ref{app:data}).  One
genuine defect of the original era was found and matters at scale: the
colour-contraction tables in both codebases were silently capped, which
above certain basis sizes corrupted matrix elements without warning; both
paths now raise instead.

Three structural observations take the reachable resolution from the
original $2K = 16$--$24$ to $2K \simeq 100$ on a laptop-class machine.
First, the inner-product (norm) matrix is \emph{exactly} block diagonal in
parton configuration --- verified by brute-force colour enumeration, not
assumed --- so orthonormalization never needs the dense $n^2$ object that
was the binding memory constraint.  Second, the Hamiltonian in the weeded
basis is sparse and gets sparser with $K$, so an iterative eigensolver pays
off once the norm is blocked.  Third, with a blockwise basis the free part
projects to an exact diagonal, eliminating the last dense intermediate.
None of these approximations the physics: the sparse path is gated to
reproduce the dense reference to solver precision, and the standard
Hamiltonian path remains bit-identical to what the code has always
produced.

Resolution is bounded above deliberately.  At $2K = 100$ the solver
reproduces its recorded reference values in seconds; past it, the colour
path returns spurious negative eigenvalues below an otherwise intact
spectrum, a failure mode that raises rather than returns in the released
code.  Requests with the wrong momentum parity --- mesons need even $2K$,
baryons $2K \equiv N \bmod 2$ --- likewise raise instead of silently
returning an empty spectrum, a trap that in our experience produces
convincing-looking ``fast runs'' otherwise.  Where the 1990 paper noted its
largest runs took ``one to three minutes of CPU on an IBM 3081,'' the
corresponding statement here is that the full table of extrapolated masses,
some eight hundred solves, rebuilds from nothing in about six minutes on
one idle workstation --- and that the binding constraint is no longer time
or memory but the endpoint physics of Sec.~\ref{sec:budget}.

% \itodo{Frozen-macro pass: replace the timing/reach statements with
% \textbackslash num macros at the v1 freeze (G1).}
